% Part: proof-theory
% Chapter: sequent-calculus
% Section: quantifiers

\documentclass[../../../include/open-logic-section]{subfiles}

\begin{document}

\olfileid{pt}{seq}{qua}

\olsection{নিয়মিত প্রমাণ ও প্রতিস্থাপন}

\LeftR{\lexists} ও~\RightR{\lforall}
বিধিতে আইগেনচল-শর্ত থাকায়
পরিমাণসূচকের বিধিগুলি কিছু
জটিলতা আনে। এই অনুমানগুলিতে
!!{constant}~$a$ আর কোথাও
ব্যবহার না করলে বেশির ভাগ
জটিলতা সামলানো যায়।
তাই একটি সংজ্ঞা দিই।

\begin{defn}
\Log{G1c} বা \Log{G3c}-র
!!^a{proof}~$\pi$ \emph{নিয়মিত},
যদি
\begin{enumerate}
  \item কোনো আইগেনচল~$a$
  একাধিক \LeftR{\lexists}
  বা~\RightR{\lforall}
  অনুমানের আইগেনচল
  হিসেবে ব্যবহৃত না হয়; এবং
  \item $a$ যদি কোনো
  \LeftR{\lexists}
  বা~\RightR{\lforall}
  অনুমানের আইগেনচল হয়,
  তবে~$\pi$-তে শুধু সেই
  অনুমানের অব্যবহিত
  উপপ্রমাণের ভিতরেই ঘটে।
\end{enumerate}
\end{defn}

\begin{lem}\ollabel{lem:subst-regular}
ধরি, $\pi$ নিয়মিত,
তার অন্তিম সিকোয়েন্ট
$\Gamma \Sequent \Delta$;
!!a{constant}~$c$ প্রমাণ~$\pi$-তে
আইগেনচল হিসেবে ব্যবহৃত নয়;
আর পদ~$t$-তে $\pi$-র
কোনো আইগেনচল নেই।
তাহলে $\pi$-তে $c$-র
প্রতিটি সংঘটন~$t$ দিয়ে
বদলে পাওয়া
$\Subst{\pi}{t}{c}$
নিয়মিত !!{proof}।
$\Subst{\pi}{t}{c}$-র
অন্তিম সিকোয়েন্ট
$\Subst{\Gamma}{t}{c}
\Sequent \Subst{\Delta}{t}{c}$।
এখানে $\Subst{\Gamma}{t}{c}$
হলো মূল বাঁ-প্রসঙ্গের প্রতিটি
সূত্রে ওই প্রতিস্থাপন করে
পাওয়া বহুসেট; প্রতিটি
সূত্রের পুনরাবৃত্তি-সংখ্যা
অক্ষত থাকে।
\end{lem}
% BN-SRC-643 TARGET-CORRECTION-BEGIN
\par\smallskip\noindent\textnormal{\small\textbf{উৎসসংশোধন (BN-SRC-643):} উৎসে \(\Gamma\)-র প্রতিস্থাপনকে সেট-গঠন দিয়ে লেখা, কিন্তু G1c/G3c-র প্রসঙ্গ বহুসেট; একই সূত্রের পুনরাবৃত্তি হারানো যায় না। সমতাটি সরিয়ে প্রতিটি সংঘটনে প্রতিস্থাপন ও পুনরাবৃত্তি-সংখ্যা সংরক্ষণ স্পষ্ট করা হয়েছে।}
% BN-SRC-643 TARGET-CORRECTION-END

\begin{proof}
  $\pheight{\pi}$-র ওপর
  আরোহ করি। $\pheight{\pi} = 0$
  হলে প্রমাণটি কেবল একটি
  স্বতঃসিদ্ধ। তখন
  $\Subst{\pi}{t}{c}$-ও
  স্বতঃসিদ্ধ: তাতে
  !!{formula}~$!A(c)$
  থাকলে তা $!A(t)$ হয়।
  পরিচয়-স্বতঃসিদ্ধের
  দুই পাশেই একই প্রতিস্থাপন
  হয়; বাঁদিকে অসংগতি-সূত্র
  থাকার স্বতঃসিদ্ধও
  অক্ষত থাকে।

  $\pheight{\pi} > 0$ হলে
  $\pi$-র শেষ অনুমানবিধি
  অনুযায়ী ক্ষেত্রভাগ করি।
  \Log{G1c}-র গঠনগত
  বিধি ও বচনীয় সংযোজকের
  যৌক্তিক বিধিগুলির ক্ষেত্র
  সহজ; অনুশীলন হিসেবে থাক।
  এখানে পরিমাণসূচকের
  বিধিতে শেষ হওয়া
  $\pi$-র ক্ষেত্রগুলি দেখি।
  \begin{enumerate}
    \item শেষ অনুমান
    $\RightR{\lforall}$
    হলে $\pi$-র রূপ:
    \[
    \AxiomC{}
    \RightLabel{$\pi'$}
    \Deduce$\Gamma(c) \fCenter \Delta'(c), !A(a,c)$
    \RightLabel{\RightR{\lforall}}
    \UnaryInf$\Gamma(c) \fCenter \Delta'(c), \lforall[x][!A(x,c)]$
    \DisplayProof
    \]
    আরোহের অনুমানে
    $\Subst{\pi'}{t}{c}$
    হলো $\Gamma(t) \fCenter
    \Delta'(t), !A(a,t)$-র
    নিয়মিত !!{proof}।
    $a$ প্রমাণ~$\pi$-র
    আইগেনচল বলে $a$
    পদ~$t$-তে ঘটে না।
    অতএব $\Subst{\pi}{t}{c}$-র
    শেষ অনুমানটি,
    \[
    \AxiomC{}
    \RightLabel{$\Subst{\pi'}{t}{c}$}
    \Deduce$\Gamma(t) \fCenter \Delta'(t), !A(a,t)$
    \RightLabel{\RightR{\lforall}}
    \UnaryInf$\Gamma(t) \fCenter \Delta'(t), \lforall[x][!A(x,t)]$
    \DisplayProof
    \]
    সঠিক: সিদ্ধান্তে~$a$
    ঘটে না।
    \item \Log{G3c}-তে
    শেষ অনুমান
    $\LeftR{\lforall}$
    হলে $\pi$-র রূপ:
    \[
    \AxiomC{}
    \RightLabel{$\pi'$}
    \Deduce$!A(s(c),c), \lforall[x][!A(x,c)], \Gamma(c) \fCenter \Delta'(c)$
    \RightLabel{\LeftR{\lforall}}
    \UnaryInf$\lforall[x][!A(x,c)], \Gamma(c) \fCenter \Delta'(c)$
    \DisplayProof
    \]
    আরোহের অনুমানে
    $\Subst{\pi'}{t}{c}$
    হলো $!A(s(t),t),
    \lforall[x][!A(x,t)],
    \Gamma(t) \fCenter
    \Delta'(t)$-র নিয়মিত
    !!{proof}।
    $\Subst{\pi}{t}{c}$-র
    শেষ অনুমানটি,
    \[
    \AxiomC{}
    \RightLabel{$\Subst{\pi'}{t}{c}$}
    \Deduce$!A(s(t),t), \lforall[x][!A(x,t)], \Gamma(t) \fCenter \Delta'(t)$
    \RightLabel{\LeftR{\lforall}}
    \UnaryInf$\lforall[x][!A(x,t)], \Gamma(t) \fCenter \Delta'(t)$
    \DisplayProof
    \]
    সঠিক। \Log{G1c}-তে
    \LeftR{\lforall}-এর
    ক্ষেত্রও অনুরূপ,
    তবে কিছুটা সরল।
    \item শেষ অনুমান
    \RightR{\lexists}
    বা \LeftR{\lexists}
    হলে: অনুশীলন।
  \end{enumerate}
  প্রতিটি ক্ষেত্রে নিয়মিত
  !!{proof}-এর শেষে
  (দুই পূর্বধারণার বিধিতে
  দুটি নিয়মিত !!{proof}s-এর
  শেষে) একটি সিকোয়েন্ট
  যোগ করেছি। নতুন অন্তিম
  সিকোয়েন্টে $\pi$-র
  অন্তিম সিকোয়েন্ট থেকে
  শুধু $c$-র স্থলে পদ~$t$
  বসেছে। $t$-তে $\pi$-র
  কোনো আইগেনচল নেই;
  তাই $\pi$ ও
  $\Subst{\pi}{t}{c}$-র
  আইগেনচলগুলি একই থাকে।
  নিয়মিত মূল প্রমাণের অন্তিম
  সিকোয়েন্টে সেগুলি
  ঘটে না, আর
  $\Subst{\pi}{t}{c}$-তেও
  নতুন করে আসে না।
  ফলে $\Subst{\pi}{t}{c}$
  নিয়মিত।
\end{proof}
% BN-SRC-644 TARGET-CORRECTION-BEGIN
\par\smallskip\noindent\textnormal{\small\textbf{উৎসসংশোধন (BN-SRC-644):} G3c-র সার্বিক-বাঁ প্রতিস্থাপনের দুটি বৃক্ষেই উৎসে ভুল \(\RightR{\lforall}\) লেবেল; পূর্বধারণা ও সিদ্ধান্ত অনুযায়ী উভয়টি \(\LeftR{\lforall}\) করা হয়েছে।}
% BN-SRC-644 TARGET-CORRECTION-END
% BN-SRC-645 TARGET-CORRECTION-BEGIN
\par\smallskip\noindent\textnormal{\small\textbf{উৎসসংশোধন (BN-SRC-645):} সহায়ক উপপাদ্যের শেষ ইংরেজি বাক্যে “and Since” দিয়ে দুটি অসম্পূর্ণ খণ্ড জুড়েছে। পূর্বধারণা অনুযায়ী \(t\)-তে পুরোনো আইগেনচল নেই, তাই প্রতিস্থাপিত প্রমাণের আইগেনচল-সমষ্টি একই ও অন্তিম সিকোয়েন্টে তারা অনুপস্থিত—এই সম্পূর্ণ যুক্তি লেখা হয়েছে।}
% BN-SRC-645 TARGET-CORRECTION-END

\begin{prop}
$\pi$ যদি \Log{G1c}
বা \Log{G3c}-তে
!!a{proof} হয়,
তাহলে একই গঠনবিশিষ্ট
(বিশেষত, একই !!{height}-এর)
নিয়মিত !!{proof}~$\pi'$
আছে, যা $\pi$ থেকে
শুধু আইগেনচলগুলি
প্রতিস্থাপন করায় আলাদা।
\end{prop}

\begin{proof}
শুধু এই প্রমাণের জন্য
$\pi$-র কোনো আইগেনচল-অনুমানকে
\emph{পরিচ্ছন্ন} বলি, যদি
তার আইগেনচল অন্য কোনো
অনুমানের আইগেনচল না হয়
এবং ওই অনুমানের অব্যবহিত
উপ-!!{proof}-র বাইরে
$\pi$-তে কোথাও না ঘটে;
অন্যথায় তাকে
\emph{অপরিচ্ছন্ন} বলি।
$\pi$-তে অপরিচ্ছন্ন
আইগেনচল-অনুমানের সংখ্যা~$n$।
$n$-র ওপর আরোহে
দেখাই, $\pi$-কে চাওয়া
নিয়মিত !!{proof}~$\pi'$-তে
বদলানো যায়।
$n=0$ হলে সব আইগেনচল-অনুমান
পরিচ্ছন্ন; তাই $\pi$
নিয়মিত এবং আর কিছু
প্রমাণের নেই। অন্যথায়
সর্বোচ্চে থাকা কোনো
\emph{অপরিচ্ছন্ন}
আইগেনচল-অনুমান বাছি,
ধরি \RightR{\lforall}।
তাতে শেষ হওয়া
$\pi$-র উপ-!!{proof}~$\pi_1$-এর
রূপ:
    \[
    \AxiomC{}
    \RightLabel{$\pi_2$}
    \Deduce$\Gamma \fCenter \Delta, !A(c)$
    \RightLabel{\RightR{\lforall}}
    \UnaryInf$\Gamma \fCenter \Delta, \lforall[x][!A(x)]$
    \DisplayProof
    \]
$c$ আইগেনচল বলে
$\Gamma$, $\Delta$ কিংবা
$\lforall[x][!A(x)]$-এ
ঘটে না। নির্বাচিত অনুমানের
ওপরে থাকা আইগেনচল-অনুমানগুলি
পরিচ্ছন্ন; তাই
উপপ্রমাণ~$\pi_2$
নিয়মিত, যদিও তাতে
আইগেনচল-অনুমান থাকতে
পারে। উপরন্তু ওই
আইগেনচল সেখানে অন্য কোনো
অনুমানের আইগেনচল নয়:
তা হলে সেটিও অপরিচ্ছন্ন
হতো। $\pi$-তে
ঘটে না এমন
!!a{constant}~$a$ নিই।
\olref{lem:subst-regular}
অনুসারে
$\Subst{\pi_2}{a}{c}$
হলো $\Gamma \fCenter
\Delta, !A(a)$-র
নিয়মিত প্রমাণ; এতে
$\RightR{\lforall}$
প্রয়োগ করা যায়।
$\pi$-তে $\pi_1$-র
বদলে প্রমাণ~$\pi_1'$
বসালে,
    \[
    \AxiomC{}
    \RightLabel{$\Subst{\pi_2}{a}{c}$}
    \Deduce$\Gamma \fCenter \Delta, !A(a)$
    \RightLabel{\RightR{\lforall}}
    \UnaryInf$\Gamma \fCenter \Delta, \lforall[x][!A(x)]$
    \DisplayProof
    \]
একটি অপরিচ্ছন্ন
আইগেনচল-অনুমান
পরিচ্ছন্ন অনুমানে
বদলেছে; স্থানীয়
উপপ্রমাণে শুধু
আইগেনচল $c$-র
জায়গায়~$a$ এসেছে।
মূল $\pi$-র অন্য কোনো
অনুমান নতুন করে অপরিচ্ছন্ন
হয় না। ফলে নতুন
প্রমাণে অপরিচ্ছন্ন
অনুমান সর্বাধিক
$n-1$টি; আরোহের
অনুমানে ফল মেলে।
\end{proof}
% BN-SRC-646 TARGET-CORRECTION-BEGIN
\par\smallskip\noindent\textnormal{\small\textbf{উৎসসংশোধন (BN-SRC-646):} \(n>0\) হলে উৎস যেকোনো সর্বোচ্চ আইগেনচল-অনুমান বেছে নেয়; সেটি পরিচ্ছন্নও হতে পারে। সর্বোচ্চ অপরিচ্ছন্ন অনুমান বেছে তার ওপরের সব অনুমান পরিচ্ছন্ন বলে \(\pi_2\) নিয়মিত দেখানো হয়েছে। একটি নাম বদলালে অন্য অপরিচ্ছন্ন অনুমানও পরিচ্ছন্ন হতে পারে, তাই নতুন সংখ্যা ঠিক \(n-1\) নয়, সর্বাধিক \(n-1\)।}
% BN-SRC-646 TARGET-CORRECTION-END

এইমাত্র প্রমাণিত প্রস্তাবটি
আমরা আলাদা করে আর
প্রয়োগ করব না; কিন্তু
বিবেচিত সব !!{proof}s-কে
নিয়মিত ধরব। নিয়মিত
না হলে আইগেনচল বদলে
সেগুলিকে নিয়মিত
করা যায়।

\end{document}
